% Part: cheat-sheets
% Chapter: quantified-modal-logics
% Section: classical-qml 

\documentclass[../../../../include/open-logic-section]{subfiles}

\begin{document}

\olfileid{chs}{qml}{cla}

\olsection{Classical Quantified Modal Logic}

The system of Classical Quantified Modal Logic, $\Log{CQML}$,
combines classical quantification, i.e. first-order logic with 
identity, with modal logic. 

Its semantics uses \emph{fixed-domain models} (also 
called \emph{constant-domain models}).

Its natural deduction system has the rules of \emph{first-order 
logic with identity} plus the modal rules.

\subsection{Fixed-domain models}

\begin{defn}[Fixed-domain models]
    A \emph{fixed-domain model} $\mModel M$ for the language $\Lang{L}$ 
    of quantified modal logic consists of the following elements:
    \begin{enumerate}
    \item \emph{Worlds:} a non-empty set~$\mWorlds{M}$
    \item \emph{Accessibility relation:} a binary relation~$\mAccess{M}$ on~$\mWorlds{M}$
    \item \emph{Domain:} a non-empty set, $\mDomain{M}$
\item \emph{Interpretation of !!{constant}s:} for each !!{constant}~$c$ of
  $\Lang{L}$, !!a{element} $\mAssign{c}{M} \in \mDomain{M}$
\item \emph{Interpretation of !!{predicate}s:} for each $n$-place
  !!{predicate}~$R$ of $\Lang{L}$ (other than $\eq$), a function 
  from worlds to $n$-place
  relations on the domain: $\mAssign{R}{M} \colon W \to \mDomain{M}^n$
\item \emph{Interpretation of !!{function}s:} for each $n$-place
  !!{function}~$f$ of $\Lang{L}$, an $n$-place function from worlds 
  to $n$-place functions on the domain: $\mAssign{f}{M} \colon W \to 
  (\mDomain{M}^n \to \mDomain{M})$

\end{enumerate}

We write $\mAssign{R}{M}[w]$ for the value of $\mAssign{R}{M}$ at $w$ (a $n$-place relation) and $\mAssign{f}{M}[w]$ for the value of the 
$\mAssign{f}{M}[w]$ at $w$ (a $n$-place function). 
\end{defn}

\begin{rem}[Model types]
  As in modal logic, a model $\mModel{M}$ is called a $\Ax{D}$-model, 
  $\Ax{T}$-model, etc., if its accessibility relation is serial,
  reflexive, etc., as in \olref{tab:five-types}. 
\end{rem}

\begin{table}[t]
    \begin{tabular}{| c | l | l |}
      \hline
      Model type & Property of $R$ & Characteristic schema \\
      \hline \hline
      \Ax{K} & no constraint & \\
      \hline
      \Ax{D} & \emph{serial}: $\forall u \exists v Ruv$ 
      & $\Box p \lif \Diamond p$ \\
      \hline
      \Ax{T} & \emph{reflexive}: $\forall w Rww$
      & $\Box p \lif p$ \\
      \hline
      \Ax{B} & \emph{symmetric}: $\forall u\forall v(Ruv \lif Rvu)$
      & $p \lif \Box\Diamond p$ \\
      \hline
      \Ax{4} & \emph{transitive}:
      $\forall u \forall v \forall w ((Ruv \land Rvw) \lif Ruw)$
      & $\Box p \lif \Box \Box p$ \\
      \hline
      \Ax{5} & \emph{euclidean}:
      $\forall w \forall u \forall v ((Rwu \land Rwv) \lif Ruv)$ 
      & $\Diamond p \lif \Box\Diamond p$ \\
      \hline
    \end{tabular}
    \caption{Five model types.}
    \ollabel{tab:five-types}
  \end{table} 

\subsection{Satisfaction in Fixed-Domain Models}

In quantified modal logics, satisfaction is relative
to an assignment (as in first-order logic or free logic)
\emph{and} to a world.

Sentences (formulas without free variables) are satisfied 
at a world irrespective of choice of assignment, so We
call them \emph{true} or \emph{false}, as in first-order logic.
 But their truth or falsity is still relative to world. 

\begin{explain}
!!^{variable} assignments and assignment variants are defined as in 
first-order logic. The !!{value} of terms is defined as in first-order
logic except that !!{function}s may vary depending on world. Note
that:
\begin{enumerate}
\item !!^{constant}s and (relative to a given assignment) !!{variable}s
    have world-invariant !!{value}s. They `directly' or `rigidly'
    refer to an object of the domain. Compare how `Socrates' denote the
    same person whether we're talking of what actually happens or 
    what could have happened.
\item Complex terms, which involve !!{function}s, have a world-sensitive
    !!{value}s. Compare how `the birthplace of Alice' or 
    `the number of planets' denote different places or numbers depending
    on how the world is.
\end{enumerate}
\end{explain}

\begin{defn}[Variable Assignment]
A \emph{variable assignment}~$s$ for a fixed-domain model~$\mModel{M}$ 
is a function which maps each !!{variable} to !!a{element} of~$\mDomain M$,
i.e., $s\colon \Var \to \mDomain M$.
\end{defn}

\begin{defn}[!!^{value} of Terms]
If $t$ is a term of the language~$\Lang L$, $\mModel M$ is a
fixed-domain model for~$\Lang L$, $w$ a world in $mWorlds M$, and $s$
is !!a{variable} assignment for~$\mModel M$, the
\emph{!!{value}}~$\mValue{t}{M}{w}[s]$ of $t$ at world $w$ in the is
defined as follows:
\begin{enumerate}
\item \indcase{t}{c}{$\mValue{\indfrm}{M}{w}[s] = \mAssign{\indcomplex}{M}$.}
\item \indcase{t}{x}{$\mValue{\indfrm}{M}{w}[s] = s(\indcomplex)$.}
\item \indcase{t}{\Atom{f}{t_1, \ldots, t_n}}{
\[
\mValue{\indfrm}{M}{w}[s] = \mAssign{f}{M}[w](\mValue{t_1}{M}{w}[s], \ldots,
\mValue{t_n}{M}{w}[s]).
\]}
\end{enumerate}
\end{defn}

$x$-variant assignments are defined as in first-order logic.
For instance, where $m$ is an object of the domain 
$\mDomain{M}$, and $s$ a !!{variable} assignment,
$\Subst{s}{m}{x}$ is the !!{variable} assignment that assigns 
$m$ to variable $x$ and is otherwise just like $s$.

\begin{defn}[Satisfaction in fixed-domain models]
Satisfaction of !!a{formula}~$!A$ a fixed-domain model~$\mModel{M}$
at a world $w$ in $\mWorlds{M}$ relative to !!a{variable} assignment~$s$, in symbols:
$\qmSat{M}{!A}[w][s]$, is defined recursively as follows. (We write
$\qmSat/{M}{!A}[w][s]$ to mean ``not $\qmSat{M}{!A}[w][s]$.'')
\begin{enumerate}
\tagitem{prvFalse}{%
  \indcase{!A}{\lfalse}{$\qmSat/{M}{\indfrm}[w][s]$.}}{}

\tagitem{prvTrue}{%
  \indcase{!A}{\ltrue}{$\qmSat{M}{\indfrm}[w][s]$.}}{}

\item \indcase{!A}{\Atom{R}{t_1, \dots, t_n}}{$\qmSat{M}{\indfrm}[w][s]$
  iff $\langle \mValue{t_1}{M}{w}[s], \dots, \mValue{t_n}{M}{w}[s] \rangle \in
  \mAssign{R}{M}[w]$.}

\item \indcase{!A}{\Atom{\lfrexists}{t}}{$\qmSat{M}{\indfrm}[w][s]$ iff 
  $\mValue{t}{M}{w}[s]\in \mDomain{M}$ (which is always the case).}

\item \indcase{!A}{\eq[t_1][t_2]}{$\qmSat{M}{\indfrm}[w][s]$ iff
  $\mValue{t_1}{M}{w}[s] = \mValue{t_2}{M}{w}[s]$.}

\tagitem{prvNot}{%
  \indcase{!A}{\lnot !B}{$\qmSat{M}{\indfrm}[w][s]$ iff
    $\qmSat/{M}{!B}[w][s]$.}}{}

\tagitem{prvAnd}{%
  \indcase{!A}{(!B \land !C)}{$\qmSat{M}{\indfrm}[w][s]$ iff 
  $\qmSat{M}{!B}[w][s]$ and $\qmSat{M}{!C}[w][s]$.}}{}

\tagitem{prvOr}{%
  \indcase{!A}{(!B \lor !C)}{$\qmSat{M}{\indfrm}[w][s]$ iff
    $\qmSat{M}{!B}[w][s]$ or $\qmSat{M}{!C}[w][s]$ (or both).}}{}

\tagitem{prvIf}{%
  \indcase{!A}{(!B \lif !C)}{$\qmSat{M}{\indfrm}[w][s]$ iff 
    $\qmSat/{M}{!B}[w][s]$ or $\qmSat{M}{!C}[w][s]$ (or both).}}{}

\tagitem{prvIff}{%
  \indcase{!A}{(!B \liff !C)}{$\qmSat{M}{\indfrm}[w][s]$ iff either
    both $\qmSat{M}{!B}[w][s]$ and $\qmSat{M}{!C}[w][s]$, or neither
    $\qmSat{M}{!B}[w][s]$ nor $\qmSat{M}{!C}[w][s]$.}}{}

\tagitem{prvAll}{%
  \indcase{!A}{\lforall[x][!B]}{$\qmSat{M}{\indfrm}[w][s]$ iff for
    every !!{element}~$m \in \mDomain M$,
    $\qmSat{M}{!B}[w][\Subst{s}{m}{x}]$.}}{}

\tagitem{prvEx}{%
  \indcase{!A}{\lexists[x][!B]}{$\qmSat{M}{\indfrm}[w][s]$ iff for at
  least one !!{element}~$m \in \mDomain M$,
  $\qmSat{M}{!B}[w][\Subst{s}{m}{x}]$.}}{}

\tagitem{prvBox}{%
  \indcase{!A}{\Box !B}{$\qmSat{M}{\indfrm}[w][s]$ iff
    $\qmSat{M}{!B}[w'][s]$ for all $w' \in \mWorlds{M}$ with
    $\mAccess{M}ww'$.}}{}

\tagitem{prvDiamond}{%
    \indcase{!A}{\Diamond !B}{$\qmSat{M}{\indfrm}[w][s]$ iff
      $\qmSat{M}{!B}[w'][s]$ for at least one $w' \in \mWorlds{M}$
      with $\mAccess{M}ww'$.}}{}

\end{enumerate}
\end{defn}

\begin{defn}[truth-at-w]
\Article{sentence} !!{sentence} $!A$ is \emph{true at a world} $w$
in model $\mModel M$ iff it is satisfied at $w$ relative to some assignment, 
or equivalently, iff it is satisfied at $w$ to every assignment.
We then simply write $\qmSat{M}{!A}[w]$. 
\end{defn}

\subsection{Natural Deduction for \Log{CQML}}

In Natural Deduction !!{derivation}s for classical 
quantified modal logic (\Log{CQML}), you may use
the following rules, detailed in the relevant 
sections of these cheatsheets:

\begin{enumerate}
  \item Propositional Logic. All rules of propositional logic.
  \item Classical Quantification and Identity. All rules of
  first-order logic.
  \item Existence. The definition of the existence predicate
  $\lfrexists$ (see below).
  \item Modality. All rules of modal logics. Depending on how the
    accessibility 
    relation of our models are constrained:
    \begin{enumerate}
      \item By default, \Log{CQML} uses modal rules from the basic
      modal system \Log{K}.
      \item Stronger systems are \Log{CQML+X}, for instance
      \Log{CQML+T} is classical quantified modal logic with modal rules
      from system \Log{T}. Its models are fixed-domain models where
      the accessibility relation is reflexive.
      \item Sider's "Simple Quantified Modal Logic" \Log{SQML}
      is \Log{CQML+S5}, classical quantified modal logic with 
      modal rules for system \Log{S5}. Its models are fixed-domain models where
      the accessibility relation is an equivalence relation (reflexive, 
      symmetric and transitive).
    \end{enumerate}
\end{enumerate}

\begin{defish}
    \emph{Definition of $\lfrexists$}

    \begin{center}        
    \smallskip
    \begin{tabular}{cc}
    \AxiomC{$\lexists[v][\eq[c][v]]$}
    \RightLabel{Def~$\lfrexists$}
    \UnaryInfC{$\lfrexists c$}
    \DisplayProof
    &
    \AxiomC{$\lfrexists c$}
    \RightLabel{Def~$\lfrexists$}
    \UnaryInfC{$\lexists[v][\eq[c][v]]$}
    \DisplayProof
    \end{tabular}
    \end{center}

\end{defish}

\end{document}